\documentclass[10pt,a4paper]{amsart}
\usepackage{amsmath,amssymb,amsthm,graphicx}
\usepackage{verbatim}
\newenvironment{reference}{\comment}{\endcomment}
%\newenvironment{reference}{}{}
\newenvironment{slogan}{\comment}{\endcomment}
\newenvironment{history}{\comment}{\endcomment}

% For commutative diagrams we use Xy-pic
\usepackage[all]{xy}

% We use 2cell for 2-commutative diagrams.
\xyoption{2cell}
% Theorem environments.
%
\theoremstyle{plain}
\newtheorem{theorem}[subsection]{Theorem}
\newtheorem{proposition}[subsection]{Proposition}
\newtheorem{lemma}[subsection]{Lemma}

\theoremstyle{definition}
\newtheorem{definition}[subsection]{Definition}
\newtheorem{example}[subsection]{Example}
\newtheorem{exercise}[subsection]{Exercise}
\newtheorem{situation}[subsection]{Situation}

\theoremstyle{remark}
\newtheorem{remark}[subsection]{Remark}
\newtheorem{remarks}[subsection]{Remarks}

\numberwithin{equation}{subsection}

% Macros
%
\def\lim{\mathop{\mathrm{lim}}\nolimits}
\def\colim{\mathop{\mathrm{colim}}\nolimits}
\def\Spec{\mathop{\mathrm{Spec}}}
\def\Hom{\mathop{\mathrm{Hom}}\nolimits}
\def\Ext{\mathop{\mathrm{Ext}}\nolimits}
\def\SheafHom{\mathop{\mathcal{H}\!\mathit{om}}\nolimits}
\def\SheafExt{\mathop{\mathcal{E}\!\mathit{xt}}\nolimits}
\def\Sch{\mathit{Sch}}
\def\Mor{\mathop{\mathrm{Mor}}\nolimits}
\def\Ob{\mathop{\mathrm{Ob}}\nolimits}
\def\Sh{\mathop{\mathit{Sh}}\nolimits}
\def\NL{\mathop{N\!L}\nolimits}
\def\CH{\mathop{\mathrm{CH}}\nolimits}
\def\proetale{{pro\text{-}\acute{e}tale}}
\def\etale{{\acute{e}tale}}
\def\QCoh{\mathit{QCoh}}
\def\Ker{\mathop{\mathrm{Ker}}}
\def\Im{\mathop{\mathrm{Im}}}
\def\Coker{\mathop{\mathrm{Coker}}}
\def\Coim{\mathop{\mathrm{Coim}}}

% Boxtimes
%
\DeclareMathSymbol{\boxtimes}{\mathbin}{AMSa}{"02}

%
% Macros for moduli stacks/spaces
%
\def\QCohstack{\mathcal{QC}\!\mathit{oh}}
\def\Cohstack{\mathcal{C}\!\mathit{oh}}
\def\Spacesstack{\mathcal{S}\!\mathit{paces}}
\def\Quotfunctor{\mathrm{Quot}}
\def\Hilbfunctor{\mathrm{Hilb}}
\def\Curvesstack{\mathcal{C}\!\mathit{urves}}
\def\Polarizedstack{\mathcal{P}\!\mathit{olarized}}
\def\Complexesstack{\mathcal{C}\!\mathit{omplexes}}
% \Pic is the operator that assigns to X its picard group, usage \Pic(X)
% \Picardstack_{X/B} denotes the Picard stack of X over B
% \Picardfunctor_{X/B} denotes the Picard functor of X over B
\def\Pic{\mathop{\mathrm{Pic}}\nolimits}
\def\Picardstack{\mathcal{P}\!\mathit{ic}}
\def\Picardfunctor{\mathrm{Pic}}
\def\Deformationcategory{\mathcal{D}\!\mathit{ef}}
\usepackage{cleveref}
\crefname{originalsection}{original section}{original sections}
\Crefname{originalsection}{Original section}{Original sections}
\makeatletter
\newlabel{algebra-section-spectrum-ring}{{00DY}{source}{}{sourcepointer-1}{}}
\newlabel{algebra-section-spectrum-ring@cref}{{[originalsection][0][]00DY}{[1][0][]source}{}{}{}}
\makeatother
\title{A quasi-affine scheme becomes an affine space after an affine base change}
\author{The Stacks Project Authors}
\date{stacks-project:properties.tex:local-snapshot-20261012; GNU FDL 1.2 or later}
\begin{document}
\maketitle
\noindent\textit{Editorial note.} All delivered mathematical text is the author's own, copied from the pinned Stacks Project chapter file properties.tex; no statement or proof step was rewritten, reordered or removed.\par\medskip
\noindent\textit{Editorial note.} The counted proof is the author's complete argument: the presentation of the quasi-affine scheme as the complement of finitely many equations, the construction of the affine scheme over X from those equations, and the explicit verification of the base change identification by mutually inverse isomorphisms.\par\medskip
\noindent\textit{Editorial note.} The label lemma-quasi-affine-presentation names the counted statement, and the proof environment immediately following it in the source is delivered verbatim as its proof.\par\medskip
\noindent\textit{Editorial note.} The counted statement is carried by the author's own lemma environment in the chapter source.\par\medskip
\noindent\textit{Editorial note.} Pre-open gate: python operations/agent-authoring/tools/precard\_gate.py --source properties.tex --label lemma-quasi-affine-presentation --start 2738 --end 2802 --local-id 28034 returned RESULT=OK (L28034 free, occupancy FREE, public interval clear). Candidate taken from sources/stacks-properties/candidates-halfA.json.\par\medskip
\noindent\textit{Editorial note.} The statement and proof are copied from the pinned chapter file; the cited inputs are kept as external pointers with their canonical Stacks tags.\par\medskip
\noindent\textit{Allowed standard prerequisites (editorial).}\begin{itemize}
\item Description of a quasi-affine scheme as a quasi-compact open of an affine scheme. Same chapter, the definition of a quasi-affine scheme and the description of a quasi-compact open of an affine scheme as a complement of a closed subset defined by finitely many equations. Supplies the presentation X = Spec(A) minus V(I) with I generated by finitely many elements f\_1, ..., f\_n.
\item Basic opens form a basis of the spectrum of a ring. Algebra, the section on the spectrum of a ring, where the standard opens D(f) are shown to form a basis of the topology. Supplies the covering of the quasi-compact open by finitely many standard opens, hence the finiteness of the list of equations.
\end{itemize}\medskip
\ifdefined\hypertarget\hypertarget{sourcepointer-1}{}\fi
\noindent\textit{Source reference.} Original section 00DY: sources/stacks-chow/tags.txt, source lines 515--515. Prerequisite pointer invoked at proof line(s) 2750; it names a standard fact of a sibling chapter that the counted argument uses, and every step of the counted proof built on it is retained verbatim in the delivered body.\par\medskip
\begin{lemma}
\label{lemma-quasi-affine-presentation}
Let $X$ be a quasi-affine scheme. There exists an integer $n \geq 0$,
an affine scheme $T$, and a morphism $T \to X$ such that for every
morphism $X' \to X$ with $X'$ affine the fibre product $X' \times_X T$
is isomorphic to $\mathbf{A}^n_{X'}$ over $X'$.
\end{lemma}

\begin{proof}
By definition, there exists a ring $A$ such that $X$ is isomorphic to a
quasi-compact open subscheme $U \subset \Spec(A)$. Recall that the standard
opens $D(f) \subset \Spec(A)$ form a basis for the topology, see
Algebra, Section \ref{algebra-section-spectrum-ring}. Since $U$ is
quasi-compact we can choose $f_1, \ldots, f_n \in A$ such that
$U = D(f_1) \cup \ldots \cup D(f_n)$. Thus we may assume
$X = \Spec(A) \setminus V(I)$ where $I = (f_1, \ldots, f_n)$. We set
$$
T = \Spec(A[t, x_1, \ldots, x_n]/(f_1 x_1 + \ldots + f_n x_n - 1))
$$
The structure morphism $T \to \Spec(A)$ factors through the open $X$
to give the morphism $T \to X$. If $X' = \Spec(A')$ and the morphism
$X' \to X$ corresponds to the ring map $A \to A'$, then the images
$f'_1, \ldots, f'_n \in A'$ of $f_1, \ldots, f_n$
generate the unit ideal in $A'$.
Say $1 = f'_1 a'_1 + \ldots + f'_n a'_n$.
The base change $X' \times_X T$ is the spectrum of
$A'[t, x_1, \ldots, x_n]/(f'_1 x_1 + \ldots + f'_n x_n - 1)$.
We claim the $A'$-algebra homomorphism
$$
\varphi :
A'[y_1, \ldots, y_n]
\longrightarrow
A'[t, x_1, \ldots, x_n, x_{n + 1}]/(f'_1 x_1 + \ldots + f'_n x_n - 1)
$$
sending $y_i$ to $a'_i t + x_i$ is an isomorphism. The claim finishes
the proof of the lemma. The inverse of $\varphi$ is given by the $A'$-algebra
homomorphism
$$
\psi :
A'[t, x_1, \ldots, x_n, x_{n + 1}]/(f'_1 x_1 + \ldots + f'_n x_n - 1)
\longrightarrow
A'[y_1, \ldots, y_n]
$$
sending $t$ to $-1 + f'_1 y_1 + \ldots + f'_n y_n$ and $x_i$ to
$y_i + a'_i - a'_i(f'_1 y_1 + \ldots + f'_n y_n)$ for $i = 1, \ldots, n$.
This makes sense because $\sum f'_ix_i$ is mapped to
$$
\begin{matrix}
\sum f'_i(y_i + a'_i - a'_i(\sum f'_j y_j)) =
(\sum f'_iy_i) + 1 - (\sum f'_j y_j) = 1
\end{matrix}
$$
To see the maps are mutually inverse one computes as follows:
$$
\begin{matrix}
\varphi(\psi(t) = \varphi(-1 +  \sum f'_i y_i) =
-1 + \sum f'_i (a'_i t + x_i) = t \\
\varphi(\psi(x_i)) = \varphi(y_i + a'_i - a'_i(\sum f'_j y_j)) =
a'_i t + x_i + a'_i - a'_i(\sum f'_ja'_jt + f'_jx_j) = x_i \\
\psi(\varphi(y_i)) = \psi(a'_i t + x_i) =
a'_i(-1 + \sum f'_j y_j) + y_i + a'_i - a'_i(\sum f'_j y_j) = y_i
\end{matrix}
$$
This finishes the proof.
\end{proof}

\end{document}
